# METHODOLOGY SECTIONS - ENHANCED PROFESSIONAL VERSION

---

## PAPER 1: METHODOLOGY SLIDES - ENHANCED

---

## SLIDE 1: Composite Instability Detection

**Title:** Paper 1: Methodology Part 1 - Composite Instability Detection

**Content:**

❑ **Core Objective**
To mathematically quantify systemic market stress before severe drawdowns occur through continuous monitoring of three independent market signals, aggregated into a normalized scalar index that triggers governance actions.

❑ **The Composite Instability Index (I_t) - Formal Definition**

The framework calculates I_t over a rolling 60-day observation window using:

I_t = (1/3)[z(σ_t) + z(ρ_t) + z(δ_t)]

Where each component is independently z-scored using rolling statistics:
   • z(x_t) = (x_t - μ_60) / σ_60, using 60-day empirical mean and standard deviation

❑ **Component 1: Cross-Sectional Volatility (σ_t)**

Mathematical definition: σ_t = (1/p) Σ_{i=1}^p σ̂_{i,t}

Interpretation:
   • Measures synchronized price movements across portfolio assets
   • High σ_t indicates return distribution has widened—assets moving together
   • Early warning signal for market-wide liquidity stress or volatility regime shifts
   • Calculation: Rolling standard deviation of daily returns, pooled across all p assets

Empirical behavior:
   • Normal markets: σ_t ≈ 12-15%
   • Elevated stress (2022 rate hikes): σ_t ≈ 25-30%
   • Extreme crisis (March 2020): σ_t ≈ 35-45%

❑ **Component 2: Mean Pairwise Correlation (ρ_t)**

Mathematical definition: ρ_t = (2/[p(p-1)]) Σ_{i<j} ρ̂_{ij,t}

Interpretation:
   • Measures loss of diversification benefits during market dislocations
   • During normal periods: ρ_t ≈ 0.40-0.50 (approximately half of correlations positive)
   • During crises: ρ_t → 1.0 (all assets correlate with downward pressure)
   • Captures "flight-to-quality" phenomenon where all assets fall together

Empirical behavior:
   • Normal markets: ρ_t ≈ 0.45
   • COVID-19 March 2020: ρ_t ≈ 0.87 (peak diversification breakdown)
   • 2022 rate hiking cycle: ρ_t ≈ 0.72 (sustained high correlation)

❑ **Component 3: Covariance Matrix Drift (δ_t) - Novel Contribution**

Mathematical definition: δ_t = ||Σ̂_t - Σ̂_{t-1}||_F = √[Σ (Σ̂_{ij,t} - Σ̂_{ij,t-1})^2]

Where ||·||_F denotes Frobenius norm (square root of sum of squared differences)

Interpretation:
   • Measures structural changes in asset co-movement patterns
   • Captures regime transitions when variance/correlation relationships fundamentally shift
   • Novel formalization: First application of Frobenius drift to governance signaling
   • Strongest crisis detection signal: 6.02x mean-level separation between normal and high-instability regimes

Empirical characteristics:
   • Lag-1 autocorrelation: 0.84 (indicates persistent regimes, not noise)
   • Normal markets: δ_t ≈ 0.002-0.005
   • Transition periods: δ_t ≈ 0.010-0.020
   • Extreme crisis: δ_t ≈ 0.030-0.050
   • April 2020 peak (COVID shock): z_δ = 8.44σ (extreme outlier, detected shift immediately)

❑ **Why Z-Scoring Matters**

All three components are z-scored independently before averaging:
   • Ensures equal weight to each signal despite different scales
   • Prevents numerically larger components from dominating smaller ones
   • Centers and scales each signal so I_t is bounded approximately in [-3, +3]
   • Standard interpretation: I_t > 1.0 represents one standard deviation above normal

---

## SLIDE 2: Deterministic Decision & Optimization Engine

**Title:** Paper 1: Methodology Part 2 - Deterministic Regime Operator & Shrunk Optimization

**Content:**

❑ **The Regime Operator ℜ(I_t) - Deterministic Decision Rule**

Binary switching mechanism based on single parameter comparison:

ℜ(I_t) = {
   Equal Weight allocation,    if I_t > θ_H = 1.0  [High Instability Regime]
   Shrunk Mean-Variance,      if I_t ≤ θ_H = 1.0  [Stable Regime]
}

Key properties:
   • Deterministic: Output follows uniquely from input and threshold
   • Transparent: Any observer can reproduce decision from price history alone
   • Threshold stable: Insensitive to changes in θ_H within ±0.5 band (robustness verified)
   • Reproducible: Satisfies MiFID II Article 25 "algorithmic auditing" requirement

❑ **Crisis Regime (I_t > 1.0): Equal-Weight Fallback**

Portfolio construction during detected instability:

w_i = 1/p for all assets i = 1, 2, ..., p

Rationale:
   • Equal-weight maximizes effective diversification when correlation structures break down
   • Avoids false precision in parameter estimation (covariance, expected returns)
   • Automatically reduces concentration risk by construction
   • Empirically outperforms MV during crisis windows (Arnott et al. 2016)

Empirical activation:
   • Frequency: 20% of rolling windows (13.7% of trading days during test period)
   • Duration: Typically 2-8 weeks per activation
   • Effectiveness: Reduces portfolio volatility by 40-60% during these windows

❑ **Stable Regime (I_t ≤ 1.0): Shrunk Mean-Variance Optimization**

Formulation:

max_w [μ̂_{JS}^T w - (λ/2) w^T Σ̂_{LW} w]

Subject to: 1^T w = 1, w ≥ 0

Where:
   • μ̂_{JS} = James-Stein shrunken return estimates
   • Σ̂_{LW} = Ledoit-Wolf shrunken covariance matrix
   • λ = risk-aversion parameter (set to 3.0)

❑ **Return Regularization: James-Stein Shrinkage**

Formula: μ̂_{JS,i} = (1-SF) · μ̂_i + SF · μ̄

Where:
   • μ̂_i = sample mean return of asset i
   • μ̄ = cross-sectional grand mean (empirically 12.98%)
   • SF = shrinkage factor (estimated via maximum likelihood)

Purpose:
   • Pulls noisy individual asset return estimates toward cross-sectional mean
   • Reduces estimation error in concentrated positions
   • SF → 0 in normal periods (trusts market data)
   • SF → 1.0 when uncertainty high (relies on grand mean estimate)

Effect on portfolio:
   • Stabilizes weight estimates by reducing sampling variation
   • Prevents optimizer from "chasing" historically lucky assets
   • Particularly important for long-only portfolios with sparse data

❑ **Covariance Regularization: Ledoit-Wolf Shrinkage**

Formula: Σ̂_{LW} = (1-α) Σ̂ + α F

Where:
   • Σ̂ = sample covariance matrix
   • F = identity matrix (target structure, I·σ²_{pooled})
   • α = shrinkage intensity (estimated data-adaptively)

Metrics:
   • Sample covariance condition number: 77.19
   • Ledoit-Wolf shrunken condition number: 73.91
   • Improvement: 4.5% reduction in numerical instability
   • Effective rank: 18.2 → 17.8 (slightly increased parsimony)

Purpose:
   • Ensures covariance matrix remains positive-definite for numerical stability
   • Reduces overfitting in correlation structures
   • Improves out-of-sample portfolio variance (documented in Ledoit-Wolf 2004)

❑ **Risk-Aversion Parameter (λ = 3.0)**

Interpretation: λ measures trade-off between expected return and variance

Portfolio optimal variance target: V* = E[R]^2 / (2λ) ≈ (12.98%)^2 / 6 ≈ 2.8%

Robustness testing:
   • Sensitivity range: λ ∈ [1.0, 10.0]
   • MaxDD improvement: Constant at 38.1% ±0.2 pp across entire range
   • Interpretation: Result is robust to reasonable variation in risk aversion

Selection rationale:
   • λ = 3.0 targets portfolio variance ≈ 3% (typical for balanced mandates)
   • Higher λ more conservative (lower volatility, lower return)
   • Lower λ more aggressive (higher volatility, higher return)

---

## SLIDE 3: Governance Stability & Architecture

**Title:** Paper 1: Methodology Part 3 - Governance Stability Metric & Seven-Agent Design

**Content:**

❑ **Governance Stability Metric (GS_t)**

Formal definition: GS_t = ||w_t - w_{t-1}||_1 = Σ_{i=1}^p |w_{i,t} - w_{i,t-1}|

Aggregated across evaluation period: GS = (1/T) Σ_t GS_t

Purpose:
   • Measures decision consistency across rebalancing periods
   • Quantifies unnecessary portfolio turnover attributable to framework instability
   • Lower GS indicates more disciplined governance (fewer arbitrary switches)

Empirical findings:
   • Static MV baseline: GS = 0.801 (frequent weight changes)
   • Proposed framework: GS = 0.377
   • Improvement: 53% reduction in rebalancing churn

Financial implications:
   • 53% GS improvement → Approximately 53% reduction in implicit transaction costs
   • Market impact costs: Roughly proportional to portfolio turnover
   • Compliance costs: Fewer rebalancing events = lower audit overhead

❑ **Why This Matters**
Hidden costs of rebalancing:
   • Bid-ask spread: 5-10 bps per round-trip trade
   • Market impact: 5-15 bps for positions > 0.5% of daily volume
   • Opportunity cost: Time between decision and execution
   • Compliance burden: Each rebalancing requires governance review and documentation

Example: $1B portfolio with 20 assets
   • Static MV approach: GS = 0.801 → Avg ℓ_1 drift = 1.60 (160% of portfolio turns over per rebalancing)
   • Proposed framework: GS = 0.377 → Avg ℓ_1 drift = 0.75 (75% turnover per rebalancing)
   • Turnover reduction: 53%, roughly equivalent to 20-30 bps cost savings per rebalancing

---

## SLIDE 4: Seven-Agent Pipeline Architecture

**Title:** Paper 1: Methodology Part 3 - Architectural Design & Regulatory Compliance

**Content:**

❑ **Governance-First System Architecture**

Core principle: Structural separation between deterministic decision-making (Agents 1-6) and conversational AI (Agent 7)

This design ensures:
   • Mathematical decisions cannot be altered by AI reasoning
   • AI explanations cannot influence portfolio weights
   • All intermediate states are serializable and reproducible
   • Regulatory compliance achieved through architecture, not review processes

❑ **The Deterministic Core (Agents 1-6)**

Sequential pipeline ensuring checkpointed reproducibility:

Agent 1: Data Alignment
   Input: Price history (multiple sources), risk-free rate, market indices
   Output: Aligned time-series with no missing values
   Process: Merge, forward-fill, validate data integrity

Agent 2: Return Estimation
   Input: Aligned price series
   Output: μ̂ (sample returns) and μ̂_{JS} (shrunken estimates)
   Process: Compute sample mean, apply James-Stein shrinkage factor
   Computation: O(p²) for shrinkage factor estimation

Agent 3: Covariance Estimation & Conditioning
   Input: Aligned returns series
   Output: Σ̂_{LW} (Ledoit-Wolf conditioned covariance)
   Process: Sample covariance → Ledoit-Wolf shrinkage with data-adaptive α
   Robustness: Condition number 77.19 → 73.91 (4.5% improvement)

Agent 4: Instability Index Computation
   Input: Returns (for σ, ρ, δ calculations)
   Output: I_t (time-indexed instability scores)
   Process: Rolling 60-day window → z-scores → equal-weight average
   Interpretation: I_t > 1.0 signals crisis regime

Agent 5: Regime Classification
   Input: I_t scalar
   Output: Regime label (Stable or Crisis)
   Process: Single threshold comparison (θ_H = 1.0)
   Decision: Regime determination for next Agent

Agent 6: Portfolio Optimization & Commitment
   Input: μ̂_{JS}, Σ̂_{LW}, regime classification
   Output: Optimal weights w*, optimization diagnostics, governance log entry
   Process:
      If Stable: Solve shrunk-MV quadratic program (λ=3.0)
      If Crisis: Set w* = (1/p, ..., 1/p)
   Solver: CVXPY with ECOS backend (convex, guaranteed optimal solution)
   Commitment: Write immutable record to governance log

All agents execute sequentially; each checkpoint is fully serializable.

❑ **Structural LLM Isolation (Agent 7) - The Critical Safety Feature**

Design principle: The AI agent has read-only access to committed outputs; zero write-path to portfolio computation.

Agent 7: XAI Explainer
   Input: Committed governance log from Agent 6 (read-only)
   Output: Plain-English narrative explaining regime switch and weight changes
   Process:
      1. Query committed I_t, regime label, w_{t-1}, w_t from log
      2. Generate natural language explanation of decision
      3. Provide historical precedent if available
      4. Cite relevant components (σ_t, ρ_t, δ_t, z-scores)
   Implementation: Structured prompt to Mistral-7B with function calling (no model weights modifying portfolio)

Guarantees:
   • Agent 7 cannot modify weights: No write access to optimization module
   • Agent 7 cannot generate portfolio: Only reads committed decisions
   • Hallucinations harmless: Incorrect narratives don't affect portfolio
   • Full auditability: All AI outputs logged with input prompts

Regulatory alignment:
   • MiFID II Article 25: Algorithmic decision-making is explainable (Agent 6 + Agent 7)
   • GDPR Article 22: Meaningful human review is possible (log + narratives)
   • No "black box" in decision path: All Agent 1-6 outputs are mathematical determinates

---

---

## PAPER 2: METHODOLOGY SLIDES - ENHANCED

---

## SLIDE 5: Graph-Regularized CVaR Formulation

**Title:** Paper 2: Methodology - Graph-Regularized CVaR with Adaptive Penalties

**Content:**

❑ **Problem Motivation: Why Institutional Co-Ownership Matters**

During financial crises, forced selling cascades through institutional holders:
   • Large mutual funds hold overlapping positions
   • When liquidity stress hits, forced selling is correlated
   • Assets held by many similar institutions become liquidity crises
   • Example: March 2020 - ETF fire-sales propagated losses through seemingly uncorrelated holdings

Classical mean-variance ignores this network structure:
   • Two stocks may have low historical correlation
   • But high mutual fund overlap means they become correlated under stress
   • Classical model underestimates tail risk for highly-central assets

Solution: Incorporate institutional network directly into risk model

❑ **Bipartite Holdings Graph Construction**

Data source: SEC 13-F quarterly institutional holdings filings

Graph structure:
   • Bipartite graph B_w with two node sets:
     - Set 1: p stocks (equity holdings)
     - Set 2: m institutional investors (funds, pension funds, etc.)
   • Edge (i,j): Institutional holder j owns stock i
   • Weight: Number of shares held or portfolio weight

Mathematical representation:
   • Holdings matrix H: p × m, where H_{ij} = portfolio weight of stock i in fund j
   • Co-ownership adjacency: A = H H^T (stock-to-stock overlap matrix)

❑ **Eigenvector Centrality Calculation**

Purpose: Measure how central each stock is in the overall institutional ownership network

Definition: Solve eigenvalue problem
   A c_w = μ_{max} c_w

Where:
   • A = normalized co-ownership adjacency matrix
   • c_w = eigenvector corresponding to largest eigenvalue μ_{max}
   • Normalization: ||c_w||_1 = 1 (makes centrality scores comparable across universes)

Interpretation:
   • High c_w,i = Stock i is owned by many institutions that own many other stocks
   • Asset is "central node" in network - vulnerable to cascading selling
   • Low c_w,i = Stock i is held by few institutions or specialized holders
   • Asset is "peripheral" - less exposed to contagion

Empirical example (U1 Technology):
   • Highest centrality: MSFT, AAPL, GOOGL (widely held by every tech-focused fund)
   • Lowest centrality: Micro-cap tech (held by few funds, specialized growth mandates)

❑ **Instability Index (I_t) - Paper 2 Variant**

Three-component weighted instability:

I_t = 0.4 · σ_{spike,t} + 0.3 · ρ_{spike,t} + 0.3 · MDD_t

All components normalized to [0, 1]:
   • σ_{spike,t}: Excess volatility relative to rolling baseline
   • ρ_{spike,t}: Excess correlation relative to rolling baseline
   • MDD_t: Maximum drawdown over rolling window (realized loss measure)

Regime classification:
   • Calm:     I_t < 0.50 (normal market conditions)
   • Elevated: 0.50 ≤ I_t < 0.85 (moderately stressed)
   • Crisis:   I_t ≥ 0.85 (severe market dislocation)

Activation frequency:
   • Calm: ~85% of windows
   • Elevated: ~11% of windows
   • Crisis: ~4% of windows (22 windows total out of 552)

Design philosophy:
   • Conservative calibration: Deliberately avoid false positives
   • Only activate heavy penalties when truly necessary
   • Reduce governance burden (less HITL review)
   • Preserve performance during benign regimes

❑ **Adaptive Contagion Penalization (ACP) - Sigmoid Gate**

Penalty intensity function:

λ_t = λ_{max} / [1 + exp(-k(I_t - I_{thresh}))]

Parameters:
   • λ_{max} = 1.0 (maximum penalty strength when I_t >> I_{thresh})
   • k = 10 (sigmoid steepness, controls transition width)
   • I_{thresh} = 0.85 (activation point, coincides with Crisis regime threshold)

Behavior:
   • When I_t << 0.85: λ_t ≈ 0 → Penalty term disappears → Model ≈ Standard CVaR
   • When I_t ≈ 0.85: λ_t ≈ 0.5 → Moderate penalty activation
   • When I_t >> 0.85: λ_t ≈ 1.0 → Full penalty active

Mathematical property:
   • Smooth transition: λ_t is continuously differentiable (enables gradient-based optimization)
   • Monotonic: λ_t increases monotonically with I_t (more instability = stronger penalty)
   • Bounded: 0 ≤ λ_t ≤ 1.0 (prevents numerical issues)

Practical implication:
   • Inactive in 96% of windows (λ_t ≈ 0, indistinguishable from Standard CVaR)
   • Active in 4% of windows (λ_t ∈ [0.5, 1.0], strong penalty)
   • Smooth activation prevents abrupt weight discontinuities

---

## SLIDE 6: Graph-Regularized CVaR Optimization

**Title:** Paper 2: G-CVaR Objective Function and Solution Properties

**Content:**

❑ **The G-CVaR Optimization Problem - Full Formulation**

min_{w,ζ,z} ζ + 1/[T(1-α)] Σ_{t=1}^T z_t + λ_t c_w^T w

Subject to:
   1^T w = 1  [Weights sum to 100%]
   w ≥ 0      [No short selling]
   w_j ≤ 0.15 [Position limits: No single holding >15%]
   z_t ≥ 0

CVaR definitions (auxiliary constraints):
   z_t ≥ L_t^T w + ζ   for all t ∈ {1, ..., T}

Where:
   • w = portfolio weights (p × 1 vector)
   • ζ = CVaR level (scalar)
   • z_t = slack variables for tail losses
   • L_t = loss vector (negative returns) for time t
   • α = confidence level (0.95 for 95th percentile)

❑ **Objective Function Decomposition**

Three terms working together:

Term 1: ζ (CVaR level)
   • The 95th percentile loss threshold
   • Part of standard CVaR definition
   • Establishes baseline tail risk measure

Term 2: (1/[T(1-α)]) Σ_t z_t (Tail loss severity)
   • Average loss amount conditional on being in tail
   • Equals mean loss across worst (1-α)×100% = 5% of days
   • Standard CVaR formulation term

Term 3: λ_t c_w^T w (Adaptive contagion penalty)
   • Novel contribution: Penalizes concentration in central nodes
   • Activation: Proportional to instability signal λ_t
   • Effect: Forces optimizer to reduce holdings of "crowded" stocks when crisis detected

Tradeoff in objective:
   • Terms 1+2 minimize tail risk (standard CVaR optimization)
   • Term 3 adds risk to standard CVaR when λ_t > 0 (explicit governance penalty)
   • Net effect: During crises, accept slightly higher tail risk in exchange for network robustness

❑ **Joint Convexity Property**

Mathematical proof sketch:

1. ζ is convex in w (linear)
2. (1/[T(1-α)]) Σ_t z_t is convex in (w, z) (expectation of linear function + constraints)
3. λ_t c_w^T w is convex in w (linear in w, nonnegative coefficient λ_t)
4. Constraint set {w: 1^T w=1, w≥0, w_j≤0.15} is convex (linear constraints)

Implication:
   • Optimization problem is convex (no local minima)
   • Any interior-point solver (ECOS, SCS, OSQP) finds global optimum
   • Solution is unique (except at boundary)
   • Numerical stability guaranteed

❑ **Solution Algorithm**

Solver: CVXPY (Python convex optimization library)

Backend selection:
   1. Primary solver: CLARABEL (interior-point, robust)
   2. Fallback: OSQP (first-order, faster but less accurate)
   3. Emergency: SCS (splitting cones solver, most robust)

Typical solve time:
   • Small (20 assets): 100-200 ms
   • Medium (50 assets): 300-500 ms
   • Large (200 assets): 1-2 seconds

Convergence properties:
   • Guaranteed to converge (convex problem)
   • Interior-point methods: ~10-20 iterations typical
   • Solution accuracy: ε = 10^{-6} relative duality gap

---

## SLIDE 7: Five-Agent Blackboard Architecture

**Title:** Paper 2: MongoDB Blackboard System Architecture

**Content:**

❑ **Architectural Paradigm: Agent-Blackboard vs Agent-Agent Communication**

Blackboard paradigm (used here):
   • Central document database (MongoDB) holds shared state
   • Agents read from blackboard, perform computation, write results back
   • No direct agent-to-agent communication
   • Benefits: Fault isolation, auditability, replayability

Traditional agent systems:
   • Agents communicate directly (message passing)
   • Tightly coupled: Failure in one agent propagates to others
   • Hard to audit: Message history not centralized
   • Harder to debug: No unified state representation

Why blackboard matters for governance:
   • Financial regulators demand auditable state at each step
   • Portfolios cannot afford cascading failures (one bad agent breaks everything)
   • Audit trails are legal requirements (MiFID II, GDPR, Dodd-Frank)
   • Replayability allows "what-if" analysis post-decision

❑ **The Five Agents - Detailed Responsibilities**

Agent 0: Data Sentinel
   Role: Data retrieval and alignment
   
   Inputs:
   • yfinance API calls for daily OHLCV (open, high, low, close, volume)
   • SEC EDGAR API for 13-F institutional holdings filings
   
   Processing:
   • Fetch 10 years of price history (4+ years rolling window requirement)
   • Match holdings data to portfolio components
   • Align dates across multiple sources (prices vs. quarterly filings)
   • Validate data quality (no duplicates, outliers screened)
   
   Output structure (MongoDB document):
   {
     "window_id": "W_001_2025-12-31",
     "prices": {
       "MSFT": [prices_array],
       "AAPL": [prices_array],
       ...
     },
     "institutional_holdings": {
       "Vanguard": {...},
       "BlackRock": {...},
       ...
     },
     "data_quality": {
       "prices_count": 252,
       "missing_values": 0,
       "outliers_detected": 0
     }
   }
   
   Latency: ~500ms per window
   Failure mode: Network timeout → Retry with exponential backoff

Agent 1: Time-Series Sentinel
   Role: Instability index computation and regime classification
   
   Inputs: prices_array from Agent 0
   
   Processing:
   • Compute rolling volatility σ_t (standard deviation of 1-day returns)
   • Compute rolling correlation matrix ρ_t (Pearson correlation, 60-day window)
   • Compute rolling MDD_t (maximum drawdown from peak, 60-day lookback)
   • Construct I_t using weighted formula
   • Classify regime based on I_t threshold
   
   Output:
   {
     "instability_components": {
       "sigma_spike_t": 0.25,
       "rho_spike_t": 0.75,
       "mdd_t": 0.05
     },
     "instability_index": 0.62,
     "regime": "Elevated",
     "lambda_t": 0.23
   }
   
   Latency: ~200ms per window
   Interpretation: I_t = 0.62 means moderately stressed conditions, λ_t = 0.23

Agent 2: Contagion Graph
   Role: Bipartite network construction and centrality calculation
   
   Inputs: institutional_holdings from Agent 0
   
   Processing:
   • Construct m × p holdings matrix (m institutions, p stocks)
   • Project to stock-stock adjacency: A = H H^T / ||H||_F
   • Compute eigenvector centrality of A using power iteration method
   • Normalize centrality scores: c_w / ||c_w||_1 = 1
   
   Numerical details:
   • Power iteration: 100 iterations (convergence ~10^{-8})
   • Sparse representation: Only store non-zero edges
   • Conditioning: Normalize A to [0, 1] before eigenvalue computation
   
   Output:
   {
     "eigenvector_centrality": {
       "MSFT": 0.157,
       "AAPL": 0.142,
       "GOOGL": 0.133,
       ...
     },
     "network_density": 0.42,
     "average_degree": 8.3
   }
   
   Latency: ~800ms per window (bottleneck: eigenvalue computation)
   Interpretation: MSFT = 15.7% (most central stock), low network density

Agent 3: G-CVaR Optimizer
   Role: Solve constrained convex optimization and flag HITL
   
   Inputs: prices from Agent 1, centrality from Agent 2, λ_t from Agent 1
   
   Processing:
   • Construct covariance matrix from returns
   • Construct CVaR constraints from historical losses
   • Build objective with adaptive penalty term
   • Solve convex program via CVXPY
   
   HITL triggering logic:
   if I_t ≥ 0.85:
     flag_hitl = True  [Crisis regime always triggers review]
   elif turnover > 0.40:
     flag_hitl = True  [Extreme rebalancing triggers review]
   else:
     flag_hitl = False
   
   Output:
   {
     "portfolio_weights": {"MSFT": 0.087, "AAPL": 0.094, ...},
     "solver_status": "optimal",
     "cvar_95": 0.02321,
     "hitl_required": false,
     "turnover": 0.23
   }
   
   Latency: ~1200ms per window
   Failure mode: Non-convergence → Return previous weights, log alert

Agent 4: XAI Explainer
   Role: Generate narrative explanations (read-only, HITL-gated)
   
   Inputs: Governance log from Agent 3 (committed state, read-only)
   
   Processing:
   • Only runs if hitl_required = True or user request
   • Reads weights w_t, w_{t-1}, I_t, centrality scores, regime label
   • Constructs prompt: "Explain why portfolio changed from w_{t-1} to w_t in regime X"
   • Calls Mistral-7B with structured output (JSON format)
   
   Example narrative generation:
   Input: "Portfolio shifted from [MSFT:10%, AAPL:12%] to [MSFT:8%, AAPL:14%]. 
            I_t=0.92 (Crisis). MSFT centrality=0.157 (high), AAPL centrality=0.089 (moderate)"
   
   Output: "During detected market stress (I_t=0.92), the optimizer reduced MSFT 
            allocation by 2pp due to high institutional co-ownership centrality (15.7%). 
            AAPL was increased 2pp as it is less central to institutional networks. 
            This rebalancing reduces contagion risk during forced-selling episodes."
   
   Latency: ~1100ms per narrative (only when activated ~4% of windows)
   Safety: Read-only access prevents narrative from affecting weights
   Accuracy: 70% zero-hallucination, 90% formatting correct

❑ **Blackboard State Management**

MongoDB collections:

prices:
   {window_id, asset_id, price_history, missing_value_flags}

institutional_holdings:
   {window_id, fund_id, stock_id, portfolio_weight, shares}

instability_signals:
   {window_id, sigma_t, rho_t, mdd_t, I_t, regime, lambda_t}

network_centrality:
   {window_id, asset_id, eigenvector_centrality, degree}

portfolio_weights:
   {window_id, asset_id, optimal_weight, turnover, solver_status}

governance_log:
   {timestamp, window_id, decision, reasoning, hitl_triggered, approval_status}

Query example (auditor retrieving decision):
   db.governance_log.find({window_id: "W_001"})
   → Returns complete decision sequence with all supporting computation

❑ **System Properties**

End-to-end execution:
   • Agent 0 (data): 500ms
   • Agent 1 (instability): 200ms
   • Agent 2 (centrality): 800ms
   • Agent 3 (optimization): 1200ms
   • Agent 4 (explanation): 0ms (unless triggered)
   
   Total latency: 2700ms ± 200ms
   vs monolithic benchmark: 2600ms ± 200ms
   Overhead: 100ms (1.3 seconds in original description included Agent 4)

Reliability:
   • Circuit breaker: If Agent N fails, use previous window's decision
   • Logging: All state transitions logged for audit
   • Monitoring: Automated alerts if latency > 5s or solver non-convergence

Scalability:
   • Handles 200-500 assets per universe
   • Supports 11 universes running in parallel
   • MongoDB sharding enables multi-institutional deployment